% Learning InitechBase, Chapter 8: Working with Several Databases (AN-16).
% Every panel replayed (examples/l8-*.dot); sources cmd/flow.c SELECT,
% cmd/workarea.c (10 areas, aliases), cmd/set.c SET FILTER.
\chapter{Working with Several Databases}
\Gates{test-pubs-examples, test-interp-use, test-interp-set}
The staff database gives each employee's department as a code; the name of the
department and the floor it is on are in another database,
\file{DEPTS.DBF}.  In this lesson you open both at once, look up a department
from an employee's record, and set a filter that makes a database show only
the records you want.

\begin{ChapterCovers}
\item Work areas and \cmd{SELECT}
\item Aliases
\item Looking up a record in another database
\item Filters
\end{ChapterCovers}

\ATEHead{Preparing for This Chapter}
Start InitechBase with a fresh copy of the sample files.

\ATEHead{Work Areas}
InitechBase has ten work areas, numbered 1 to 10.  Each can hold one open
database, with its index files and its own record pointer.  \cmd{SELECT}
chooses the work area in which the next commands work, and \cmd{USE} opens a
database in the selected work area.  Work area 1 is selected when InitechBase
starts.
\begin{Steps}
\item Type \type{SELECT 1}, then \type{USE STAFF}.
\item Type \type{SELECT 2}, then \type{USE DEPTS INDEX DEPTS}.
\item Type \type{SELECT 1}, then \type{? DBF(), LASTNAME}.
\item Type \type{SELECT 2}, then \type{? DBF(), DEPTNAME}.
\end{Steps}
\DotExample{examples/l8-areas}
\fn{DBF()} names the database of the selected work area.  Each database keeps
its own current record.

A database can be open in only one work area at a time:
\DotExample{examples/l8-inuse}

\ATEHead{Aliases}
The alias of a database is the name by which \cmd{SELECT} knows it: its file
name without the extension, or a name you give with \kw{ALIAS}.
\DotExample{examples/l8-alias}

\ATEHead{Looking Up a Department}
To find the department of an employee, keep the department code of the
employee's record in a memory variable, select the departments, and seek the
code in their index.
\begin{Steps}[resume]
\item Type \type{SELECT 2}, then \type{USE DEPTS INDEX DEPTS}.
\item Type \type{SELECT 1}, then \type{USE STAFF}, then \type{GO 12}.
\item Type \type{MDEPT = DEPT}.
\item Type \type{SELECT DEPTS}, then \type{SEEK MDEPT}, then \type{? DEPTNAME,
FLOOR}.
\item Type \type{SELECT STAFF}, then \type{? LASTNAME, DEPT}.
\end{Steps}
\DotExample{examples/l8-lookup}
Carol Pruitt works in Accounting, on the second floor.  The staff database is
still on her record.

\ATEHead{Filters}
SET FILTER makes a database show only the records that meet a condition, until
the filter is removed.  \cmd{LIST}, \cmd{GO TOP} and \cmd{SKIP} pass over the
other records.
\begin{Steps}[resume]
\item Type \type{USE STAFF}.
\item Type \type{SET FILTER TO DEPT = 'SOFT'}, then \type{LIST LASTNAME,
FIRSTNAME}.
\item Type \type{GO TOP}, \type{? RECNO()}, \type{SKIP}, \type{? RECNO()}.
\item Type \type{SET FILTER TO} to remove the filter.
\end{Steps}
\DotExample{examples/l8-filter}
The condition of a filter may be as detailed as any \kw{FOR} condition:
\DotExample{examples/l8-filter-2}

\ATEHead{Where Do You Go From Here?}
You have worked with two databases and with a filter.  Looking up a department
for every employee by hand is slow; a program does it for you.  In Chapter~9 you
run your first programs.

\ATEHead{Summary}
\begin{enumerate}
\item \cmd{SELECT} \xp{n} or \cmd{SELECT} \xp{alias} chooses a work area; each
  has its own database and record pointer.
\item \cmd{USE} \ldots{} \kw{ALIAS} gives a database another name.
\item To use a value of one database in another, store it in a memory variable.
\item SET FILTER TO \xp{condition} hides the records that do not meet the
  condition; SET FILTER TO removes the filter.
\end{enumerate}
